% ============================================================================
% Appendix H: Production A/B Testing Guide
% Practical recipes for validating banditGPT against incumbent routers
% ============================================================================

\section{A/B Testing banditGPT in Production}
\label{app:ab_testing}

Deploying a bandit router is itself a decision under uncertainty: teams must validate that \textbf{banditGPT} outperforms their incumbent system before committing production traffic.
This appendix provides four progressively riskier integration patterns---from zero-risk shadow logging to full traffic interleaving---alongside concrete code, metric definitions, and statistical guidance.
Table~\ref{tab:ab_phases} summarises the phases; practitioners should advance through them sequentially, gating each transition on the success criteria defined below.

\begin{table}[htbp]
\centering
\caption{A/B testing phases for production adoption of banditGPT.}
\label{tab:ab_phases}
\small
\resizebox{\columnwidth}{!}{%
\begin{tabular}{@{}clccp{4.0cm}@{}}
\toprule
\textbf{Phase} & \textbf{Pattern} & \textbf{Traffic Risk} & \textbf{Learns Online} & \textbf{Gate to Next Phase} \\
\midrule
0 & Shadow (log-only)           & None   & \checkmark & Simulated reward $\geq$ incumbent \\
1 & Canary split (5--10\%)       & Low    & \checkmark & No SLA violations; reward $\geq$ incumbent ($p<0.05$) \\
2 & Balanced A/B (50/50)        & Medium & \checkmark & Statistically significant reward lift \\
3 & Full rollout + holdback (95/5) & High & \checkmark & Continuous monitoring; rollback on regression \\
\bottomrule
\end{tabular}%
}
\end{table}

% ────────────────────────────────────────────────────────────────────────────
% H.1: Phase 0 — Shadow Deployment
% ────────────────────────────────────────────────────────────────────────────

\subsection{Phase 0: Shadow Deployment (Log-Only)}
\label{app:ab_shadow}

In shadow mode, every production request is routed by the \emph{incumbent} system (static router, manual rules, or single-model) while banditGPT runs in parallel, recording what it \emph{would have} selected without affecting the user.
This provides a zero-risk baseline comparison and trains the bandit on live traffic features before it handles real requests.

\paragraph{Implementation.}

{\small\begin{verbatim}
import copy
from bandit_gpt import BanditRouter

# Production router (incumbent)
production_router = get_incumbent_router()

# Shadow banditGPT instance (log-only)
shadow = BanditRouter.create(
    cost_penalty=0.5,
    exploration="aggressive",
    use_corralling=True,
)

def handle_request(prompt: str) -> str:
    # --- Incumbent serves the user ---
    prod_model = production_router.route(prompt)
    response = call_llm(prod_model, prompt)

    # --- Shadow: banditGPT logs its counterfactual ---
    shadow_model, shadow_log = shadow.route(prompt)

    emit_metric("shadow.selected_model", shadow_model)
    emit_metric("shadow.predicted_utility",
                shadow_log.predicted_utility)
    emit_metric("shadow.cost_usd", shadow_log.cost_usd)

    # Feed the incumbent's observed reward to the shadow
    # (off-policy signal; the bandit still learns from it
    #  because the context was embedded identically)
    reward = evaluate_response(prompt, response)
    shadow.process_feedback(
        shadow_log.request_id, reward
    )
    return response
\end{verbatim}}

\paragraph{What to measure.}
\begin{itemize}[nosep,leftmargin=*]
    \item \textbf{Counterfactual agreement rate}: fraction of requests where the shadow router selects the same model as the incumbent.  A low rate ($<$50\%) indicates the bandit is finding a meaningfully different policy.
    \item \textbf{Simulated cost savings}: $\sum_t c_{\mathrm{incumbent}}(t) - c_{\mathrm{shadow}}(t)$, using the shadow's selected model cost.
    \item \textbf{Selection entropy}: $H = -\sum_a p_a \ln p_a$ over the shadow's model selection distribution.  Higher entropy during early exploration that decays over time confirms the bandit is converging.
    \item \textbf{Latency overhead}: time spent in \texttt{shadow.route()} (typically $<$5\,ms, dominated by embedding inference).
\end{itemize}

\paragraph{Gate criterion.}
Advance to Phase~1 when: (1)~the shadow has processed $\geq$500 requests (sufficient for LinUCB convergence; Section~4.2), (2)~simulated reward $\geq$ incumbent reward on overlapping prompts, and (3)~latency overhead is within SLA headroom.

% ────────────────────────────────────────────────────────────────────────────
% H.2: Phase 1 — Canary Split
% ────────────────────────────────────────────────────────────────────────────

\subsection{Phase 1: Canary Split (5--10\%)}
\label{app:ab_canary}

A canary deployment routes a small fraction of live traffic through banditGPT while the incumbent handles the remainder.
This limits blast radius: even if the bandit makes poor selections, $\leq$10\% of users are affected.

\paragraph{Implementation.}

{\small\begin{verbatim}
import random
from bandit_gpt.storage import CheckpointManager

router_bandit = BanditRouter.create(
    cost_penalty=0.5,
    exploration="balanced",
    use_corralling=True,
)
checkpoint = CheckpointManager()
checkpoint.load(router_bandit)  # resume shadow state

CANARY_FRACTION = 0.05  # 5% of traffic

def handle_request(prompt, user_id):
    if hash(user_id) % 100 < CANARY_FRACTION * 100:
        # --- Treatment: banditGPT ---
        model, log = router_bandit.route(prompt)
        group = "bandit"
    else:
        # --- Control: incumbent ---
        model = production_router.route(prompt)
        log = None
        group = "control"

    response = call_llm(model, prompt)
    reward = evaluate_response(prompt, response)

    # Log to experiment tracker
    emit_ab_event(
        user_id=user_id, group=group,
        model=model, reward=reward,
        cost=get_cost(model),
    )

    # Update bandit on its own decisions only
    if log is not None:
        router_bandit.process_feedback(
            log.request_id, reward
        )
    return response
\end{verbatim}}

\paragraph{Assignment strategy.}
Hash-based assignment (\texttt{hash(user\_id) \% 100}) ensures \emph{sticky} group membership: the same user always sees the same router within a session, preventing within-user variance from confounding results.  For session-level routing (e.g., chatbots), assign at the session level rather than per-request.

\paragraph{Checkpoint and rollback.}
The \texttt{CheckpointManager} enables atomic state persistence (Section~3):
\begin{itemize}[nosep,leftmargin=*]
    \item \textbf{Before canary start}: save the shadow-trained state as a known-good checkpoint.
    \item \textbf{During canary}: periodic saves (e.g., every 100 requests) via a background thread.
    \item \textbf{On regression}: reload the pre-canary checkpoint and revert to shadow mode.
\end{itemize}

\paragraph{Gate criterion.}
Advance to Phase~2 when: (1)~no SLA violations (cost ceiling, latency) in the canary group, (2)~mean reward in the bandit group $\geq$ control ($p < 0.05$, two-sided Mann-Whitney $U$), and (3)~the canary has run for $\geq$1{,}000 requests (statistical power $\geq$0.8 for a 5\% effect size at $\alpha{=}0.05$, assuming typical reward variance).

% ────────────────────────────────────────────────────────────────────────────
% H.3: Phase 2 — Balanced A/B Test
% ────────────────────────────────────────────────────────────────────────────

\subsection{Phase 2: Balanced A/B Test (50/50)}
\label{app:ab_balanced}

A balanced split maximises statistical power for detecting reward differences.  At this phase, the bandit has already demonstrated safety in the canary; the goal shifts to precise effect-size estimation.

\paragraph{Implementation.}
The code follows the canary pattern (Section~\ref{app:ab_canary}) with \texttt{CANARY\_FRACTION\,=\,0.50}.
The only architectural addition is parallel metric collection for both arms.

\paragraph{Primary metric.}
Define the \emph{cost-adjusted reward} as the objective function banditGPT optimises (Equation~1):
\begin{equation}
    R_{\lambda}(t) = r_t - \lambda \cdot \tilde{c}_t
    \label{eq:ab_metric}
\end{equation}
where $r_t$ is the observed reward and $\tilde{c}_t$ is the normalised cost.  This single scalar captures both quality and cost, avoiding the multiple-testing problem of evaluating them separately.

\paragraph{Statistical analysis.}
\begin{enumerate}[nosep,leftmargin=*]
    \item \textbf{Primary test}: Two-sided Mann-Whitney $U$ on $R_\lambda$ between groups, with Bonferroni correction if testing multiple $\lambda$ values.
    \item \textbf{Effect size}: Report Cohen's $d$ alongside $p$-values; a $d \geq 0.2$ (small effect) is practically meaningful for routing, as it compounds over millions of requests.
    \item \textbf{Confidence interval}: Bootstrap 95\% CI on the mean reward difference $\Delta \bar{R}_\lambda = \bar{R}_\lambda^{\mathrm{bandit}} - \bar{R}_\lambda^{\mathrm{control}}$.
    \item \textbf{Segmented analysis}: Break results down by prompt difficulty (using the embedding's PC1 quintiles from Section~\ref{sec:motivation}) to identify where banditGPT's adaptive routing adds the most value.
\end{enumerate}

\paragraph{Monitoring guardrails.}
Deploy automatic stopping rules to protect against sustained regression:
\begin{itemize}[nosep,leftmargin=*]
    \item \textbf{Cost guardrail}: If the bandit group's 24-hour rolling average cost exceeds the control by $>$20\%, trigger an alert and reduce to canary fraction.
    \item \textbf{Quality guardrail}: If the bandit group's 24-hour rolling reward drops below the control's by $>$1 standard deviation, pause the experiment.
    \item \textbf{Latency guardrail}: If P99 routing latency (embedding + bandit selection) exceeds 50\,ms, investigate feature-service bottlenecks before continuing.
\end{itemize}

\paragraph{Gate criterion.}
Advance to Phase~3 when: (1)~$\Delta \bar{R}_\lambda > 0$ with $p < 0.05$, (2)~no guardrail triggers for $\geq$48 consecutive hours, and (3)~segmented analysis confirms no subpopulation regression (i.e., the bandit does not sacrifice quality on hard prompts to save cost on easy ones).

% ────────────────────────────────────────────────────────────────────────────
% H.4: Phase 3 — Full Rollout with Holdback
% ────────────────────────────────────────────────────────────────────────────

\subsection{Phase 3: Full Rollout with Holdback}
\label{app:ab_rollout}

After a successful balanced A/B test, route 95\% of traffic through banditGPT while maintaining a 5\% holdback on the incumbent.
The holdback serves two purposes: (1)~continuous monitoring for long-horizon regression (distribution drift, model deprecation), and (2)~a clean counterfactual for ongoing cost-savings reporting.

\paragraph{Ongoing operations.}
\begin{itemize}[nosep,leftmargin=*]
    \item \textbf{Checkpoint cadence}: Save state every $N{=}500$ requests or every 6 hours (whichever is more frequent).  Retain the last 3 checkpoints for rollback.
    \item \textbf{Model onboarding}: When adding new models via \texttt{register\_model()} (Section~\ref{app:scenario_quality}), temporarily increase the holdback to 10\% for 48~hours to measure the onboarding impact.
    \item \textbf{Drift detection}: Monitor the Corralling expert weights ($w_{\mathrm{warmup}}$ vs.\ $w_{\mathrm{tabula\text{-}rasa}}$).  If the tabula-rasa expert's weight exceeds 0.7 for $>$24~hours, this signals significant distribution shift---the warmup priors may need refreshing.
    \item \textbf{Quarterly review}: Re-run the balanced A/B test (Phase~2) at least quarterly to confirm continued lift, especially after major model portfolio changes.
\end{itemize}

% ────────────────────────────────────────────────────────────────────────────
% H.5: Comparing banditGPT Configurations
% ────────────────────────────────────────────────────────────────────────────

\subsection{Comparing banditGPT Configurations Against Each Other}
\label{app:ab_config}

Beyond validating banditGPT against an incumbent, practitioners may wish to compare banditGPT configurations---e.g., Corralling on vs.\ off, different $\lambda$ values, or different exploration schedules.
The library's \texttt{\_\_deepcopy\_\_} support enables this by cloning a router's learned state into independent instances:

{\small\begin{verbatim}
import copy

# Baseline: current production config
router_a = BanditRouter.create(
    cost_penalty=0.5,
    use_corralling=True,
)

# Variant: more aggressive cost optimisation
router_b = copy.deepcopy(router_a)
router_b.cost_penalty = 0.8
\end{verbatim}}

Both routers share identical learned state at fork time, isolating the effect of the configuration change.
Route traffic using the same hash-based assignment from Section~\ref{app:ab_canary}, and apply the same statistical framework from Section~\ref{app:ab_balanced}.

\paragraph{Caution: independent learning paths.}
After the fork, each router's bandit state diverges because it receives feedback only from its own routing decisions.
For short experiments ($<$500 requests per arm), this divergence is negligible.
For longer experiments, interpret results as the joint effect of the configuration change \emph{and} the different exploration trajectories it induces---which is the quantity of practical interest for deployment decisions.

% ────────────────────────────────────────────────────────────────────────────
% H.6: Summary
% ────────────────────────────────────────────────────────────────────────────

\subsection{Summary}
\label{app:ab_summary}

The phased approach above minimises production risk while providing statistically rigorous evidence for adoption.
Key design principles:
\begin{enumerate}[nosep,leftmargin=*]
    \item \textbf{Start with zero risk.}  Shadow deployment (Phase~0) trains the bandit on live features at no user-facing cost.
    \item \textbf{Gate on evidence.}  Each phase transition requires a pre-defined statistical criterion, preventing premature rollout.
    \item \textbf{Always maintain a counterfactual.}  Even at full rollout, a holdback group enables continuous monitoring and regression detection.
    \item \textbf{Leverage the bandit's own infrastructure.}  \texttt{CheckpointManager} for rollback, \texttt{process\_feedback} for delayed rewards, and \texttt{get\_probabilities} for posterior diagnostics are already built into the library---no external A/B testing platform is required (though integration with one is straightforward).
\end{enumerate}
